\documentclass[12pt]{article}
\usepackage[a4paper,margin=1in]{geometry}
\usepackage{amsmath,amssymb}
\usepackage{mdframed}
\usepackage{xcolor}

\setlength{\parskip}{1em}  % Set the space between paragraphs
\setlength{\parindent}{0pt}  % Remove the indentation at the beginning of paragraphs

\mdfsetup{
topline=false,
rightline=false,
bottomline=false,
linewidth=3pt,
innerleftmargin=15pt,
innerrightmargin=0pt,
skipabove=\baselineskip,
skipabove=1.2\baselineskip,
}

\begin{document}

\title{subsectionction exponentielle}
\author{B. Ruben}
\date{\today}

\maketitle

\section{Definitions et propriétés}

\begin{mdframed}
    \textbf{Definition:} Il existe une fonction $f$ et une seule définie et dérivable sur $\mathbb{R}$ telle que:\\ pour tout réel $x$, $f'(x) = f(x)$ et $f(0) = 1$.
    \\ Cette fonction est appelé fonction exponentielle et notée exp':$x \mapsto$ exp($x$).
    \\ On notera $e^x = \text{exp}(x)$, la derivée de $x \mapsto e^{x}$ est elle même.
\end{mdframed}

\begin{mdframed}
    \textbf{propriétés algebriques:}

    $e^{x+y} = e^x \cdot e^y$ \\ $e^{x-y} = \dfrac{e^x}{e^y}$ \\ $e^{-y} = \dfrac{1}{y}$
\end{mdframed}

\subsection*{Exercices}

$e^5 = e^3 \cdot e^2 \approx 140$, $e^{-3} = \dfrac{e^0}{e^3} \approx 0.05$, $e^1 = e$

$A = e^{2x+2}$, $B = e^{5x}$, $C = e^{-3}$


4. p.192
$e^2 = 2.7^2 \approx 7$, $3e^1 + 2 \approx 3 \cdot 2.7 + 2 = 10.1$

5. $e^2 = e^5 \cdot e^{-3} \approx{ 140 \times 0.05 = 7 }$

1. p. 191

1. $f'(x) = 3e^{x} - 2$

2. $e^x\cdot4x + e^x \cdot 2x^2 + 1$\\
racines : $x_1 = -1$, $x_2 = -1+\frac{\sqrt{2}}{2}$

3. $h(x) = \dfrac{2+e^x}{6+2x}$\\
$h = \dfrac{u}{v}$ donc $h' = \dfrac{u'v-uv'}{v^2}$\\
u = $2 + e^x$\\
u' = $e^x$\\
v = $2x+6$\\
v = $2$

donc $h'(x)$ est:
$$
    \frac{e^x(2x+6) - 2(2 + e^x)}{(2x+6)^2}
$$

Ex. 2 p 191\\
2. $g(x) = (x^2 + 2x-1)e^x$\\
$g'(x) = (2x + 2)e^x + (x^2 + 2x-1)e^x = e^x(x^2 + 4x + 1)$

Ex 20\\
1. $e^{2x} = e^0$ donc $x = 0$\\
2. impossible

\section{La fonction}

$exp: \mathbb{R} \rightarrow{} ]0; +\inf[$\\
$x \mapsto e^x$\\
$f(x) = e^x$ est positive et croissante sur $\mathbb{R}$

$e^x > 0$; $e^1 = e \approx 2,718$; $e^0 = 1$; \\ $e^a = e^b \iff a = b$ (également vrai pour les inégalités)

\section{les fonctions $x \mapsto exp(kx)$}
On nomme fonction exponentielle toutes les fonctions de la forme $f(x) = e^{kx}$ avec $k \in \mathbb{R}^*_+$, ou $g(x) = e^{-kx}$

Dérivées: $f'(x) = ke^{kx}$\\
$g'(x) = -ke^{-kx}$

$ax+b = 0 \iff x = \frac{-b}{a} $

\end{document}
